% discussion/axtran_reinterpretation-DE.tex
\chapter{axtranNew: Von der Übergangsgrammatik zur Alignment-Berechnung}
\label{chap:axtran-reinterpretation}

\section*{Motivation}

Berlinish wird erst dann zur Ingenieurmaschinerie, wenn Funktionsidentitäten,
Komponentenlängen, Randbedingungen und Freiheitsgrade gemeinsam gelöst werden
können. axtranNew entstand als diese Berechnungsantwort. Es ist mehr als ein
bequemer Optimierer: Seine vorgesehene Rolle ist der
Alignment-Berechnungskern, der eine Übergangsgrammatik in verbundene,
editierbare Alignment-Struktur überführt.

Die Bezeichnung gehört zur Implementierungslinie von ufAIM. Sie bezeichnet
keinen genehmigten Begriff des Knowledge Kernel, und ihre ursprüngliche Absicht
darf nicht mit dem gegenwärtigen Implementierungsstand verwechselt werden.

\section{Ursprüngliche Absicht}

Der vorgesehene Umfang von axtranNew umfasst:
\begin{itemize}
\item das Erfüllen von Anschluss-, Pose-, Krümmungs- und Randbedingungen;
\item die ausdrückliche Behandlung fester, abgeleiteter, beschränkter,
gekoppelter und freier Größen;
\item die Komposition von Halbwellen-, Klothoiden- und
Null-Längen-Komponenten;
\item die Bestimmung zulässiger Parameterkombinationen;
\item die Kopplung oder Trennung von Krümmungs- und Überhöhungsfunktionen;
\item den Vergleich von Alternativen und die Optimierung anhand von
Ingenieurzielen;
\item elementare Bearbeitungen wie das Koppeln von Alignment-Elementen;
\item die Realisierung berechneter Kandidaten als editierbare
Alignment-Strukturen.
\end{itemize}

Diese Liste beschreibt die Ingenieurabsicht. Sie behauptet nicht, dass jeder
Punkt produktionsreif ist.

\section{Mathematische Fähigkeit}

Mathematisch erhält ein axtranNew-Problem Funktionsidentitäten,
Parametrisierungen, Randbedingungen, Ingenieurbedingungen und ausgewiesene
Freiheitsgrade. Es erzeugt einen Kandidatenparametersatz und den daraus
abgeleiteten Alignment-Zustand. Die Berlinish-Komposition ist ein zulässiges
Strukturmodell; andere dünn besetzte Übergangsfamilien können dieselben
Berechnungsprinzipien verwenden.

Die Berechnung kann numerische Konsistenz, Residuen, Anwendbarkeit unter
angegebenen Bedingungen und Zielwerte ermitteln. Sie kann weder entscheiden,
dass der Kandidat angenommen ist, noch Engineering Object Identity übertragen
oder ihr eigenes Ergebnis autoritativ machen.

\section{Beobachtete aktuelle Implementierung}

Die aktuelle Referenzimplementierung belegt einen begrenzten, aber
substanziellen Berechnungspfad:
\begin{itemize}
\item Eine validierte Registry löst benannte Übergangsfamilien, normierte
Prototypfunktionen, Halbwellen und dreiteilige Längenteilungen auf.
\item Dünn besetzte feste und Übergangselemente werden mit fortlaufender
Pose-Übertragung und Krümmungsintegration zusammengesetzt.
\item Krümmungsträger der Länge null und ausgewählte Übergangsfälle der Länge
null werden repräsentiert.
\item Position, Tangente, Pose, Krümmung und World--Track-Operationen werden
ausgewertet.
\item Native Bearbeitungen von Geraden, Bögen und Übergängen lösen eine
deterministische dünn besetzte Neuberechnung aus und erhalten die Identität
unveränderter Elemente.
\item Ausgewählte Probleme zum Wechsel des Übergangstyps verwenden
Equality-QP- und experimentelle SQP-Maschinerie.
\end{itemize}

Diese beobachtete Implementierung unterstützt Übergangsvergleich,
Rekonstruktion und ausgewählte Bearbeitungen. Sie belegt noch keinen
allgemeinen, zuverlässigen und produktionsreifen Optimierer für jeden
Anschluss, jede Überhöhungskopplung, jedes Netz oder jedes Ingenieurziel.

\section{Gegenwärtige Grenze und künftige Entwicklung}

Der weiter reichende SQP-Dienst bleibt experimentell, nicht autoritativ und
auf Vorschläge beschränkt. Zuverlässige Konvergenz, Globalisierung,
vollständige Bedingungs- und Zielbibliotheken, validierte gemeinsame
Krümmungs--Überhöhungs-Optimierung und Berechnung auf Netzebene bleiben
Entwicklungs- oder Forschungsarbeit. Ergebnisse müssen als prüfbare
Kandidaten, Diagnosen oder Deltas zurückgegeben und vom verantwortlichen
externen Arbeitsablauf angewendet werden.

\section{Von der Geometrieoptimierung zur Laufzeitzustandsoptimierung}

AIM verallgemeinert die ursprüngliche Alignment-Berechnung zum
realisierungsbewussten Entwurf von Laufzeitzustandstrajektorien.

Dies ist keine Zurückweisung von axtranNew oder klassischer Systeme. Es erklärt,
welche zusätzlichen Unterscheidungen erforderlich sind, wenn ein berechneter
Übergang Identität, Provenienz, Realization Context, Beobachtungsbezüge,
anwendbares Wissen und verantwortlichen Entscheidungsstatus bewahren muss.

\section{Verbindung zu AIM}

\begin{summary}
Berlinish stellt die verallgemeinerte funktionale Grammatik der
Eisenbahnübergänge bereit. axtranNew macht diese Grammatik rechnerisch
anwendbar. Die weitere ufAIM-Umgebung entstand, um das resultierende Wissen zu
bewahren, anzuwenden, zu untersuchen, zu erklären und zu erweitern.

Die aktuelle Implementierung realisiert wichtige Teile dieser Linie, während
ein allgemeiner produktionsreifer Optimierer und validierte gemeinsame
Krümmungs--Überhöhungs-Optimierung künftige Arbeit bleiben. Berechnung erzeugt
Kandidaten; sie schafft keine Autorität und bildet keine Ingenieurentscheidung.
\end{summary}
